% ==> Test thm style warning msg
\InputIfFileExists{zlatex-cfg.tex}{}{}
\documentclass[
  % hyper,
]{../code/zlatex}

\zlatexThmStyle{aaa}
\begin{document}
AAA
\end{document}




% ==> Make Metropolis slide theme
\InputIfFileExists{zlatex-cfg.tex}{}{}
\documentclass[
  hyper,
  layout={slide, aspect=16|9, theme=AnnArborSpruce},
]{../code/zlatex}
\zslideThemeUse[
  UL={bg=white, text=}, 
  UR={bg=white, text=},
  BL={bg=white, text=},
  BC={bg=white, text=},
  BR={bg=white, text=},  
]{AnnArborSpruce}
\ExplSyntaxOn
\dim_gset:Nn \g_zslide_status_bar_sec_H_dim {2.2em}
\dim_gset:Nn \g_zslide_status_bar_sec_B_dim {-2.2em}
\dim_gset:Nn \g_zslide_status_info_sec_B_dim {-1.55em}
\dim_gset:Nn \g_zslide_status_info_head_B_dim {-2.5em}
% navigator
\AddToHook{shipout/foreground}{
  \_zslide_status_info:nnnn {head}{0}{1}{
    \zslideNavigateBall[
      \RPITEM[red]
    ][
      \RPITEM
    ]
  }
}
\newcommand\ITEM[1][gray]{\hspace*{-4pt}\pgfornamenthan[color=#1, scale=0.06]{39}}
\newcommand\RPITEM[1][gray]{\prg_replicate:nn {10}{\ITEM[#1]}}
\ExplSyntaxOff
\usepackage{pgfornament-han}


\title{zSlide Metropolis Theme}
\author{Eureka}
\date{\today}
\begin{document}
\maketitle


\section{SECTION AAA}
AAA

\zslideNavigateBall[\ITEM[red]][\ITEM]

\newpage
% \section{SECTION BBB}
BBB 
\newpage
% \section{SECTION CCC}
CCC 
\newpage
% \section{SECTION DDD}
DDD 

\end{document}




% ==> Test zlatex working condition
\InputIfFileExists{zlatex-cfg.tex}{}{}
\documentclass[
  % hyper,
  layout={slide, aspect=16|9, theme=AnnArborSpruce},
]{../code/zlatex}
\zslideThemeUse[
  UR={text=\zLaTeX{} Slide:\ \zslideNavigateBall},
]{AnnArborSpruce}
\zslidelogo[width=4em, exclude={0, 2, 4}]{example-image-duck}

% BEGIN TEST INTERFACE
\ExplSyntaxOn
\AddToHook{shipout/background}{
  \_zslide_status_bar:nnnn {BC}{(.5\paperwidth, -.5\paperheight)}{.2}{80pt}
}
\AddToHook{shipout/foreground}{
  \_zslide_status_info:nnnn {foot}{.66}{.75}{AAA\zslideNavigateBall[\textcolor{red}{XX}][\textcolor{gray}{XX}]}
}
\ExplSyntaxOff
% END TEST INTERFACE

\title{TEST}  
\author{Eureka}
\date{\today}
\begin{document}
\maketitle


\section{FIRST SEC}
Hello world, 

\newpage

aaa 

\newpage

bbb

\section{SECOND SEC}
this is me.

\section{THIRD SEC}
HAHA.

\end{document}



% ==> l3doc function env
\keys_define:nn { l3doc/function }
{
  TF .value_forbidden:n = true ,
  TF .code:n =
    {
      \bool_set_true:N \l__codedoc_macro_TF_bool
    } ,
  EXP .value_forbidden:n = true ,
  EXP .code:n =
    {
      \bool_set_true:N \l__codedoc_macro_EXP_bool
      \bool_set_false:N \l__codedoc_macro_rEXP_bool
    } ,
  rEXP .value_forbidden:n = true ,
  rEXP .code:n =
    {
      \bool_set_false:N \l__codedoc_macro_EXP_bool
      \bool_set_true:N \l__codedoc_macro_rEXP_bool
    } ,
  pTF .value_forbidden:n = true ,
  pTF .code:n =
    {
      \bool_set_true:N \l__codedoc_macro_pTF_bool
      \bool_set_true:N \l__codedoc_macro_TF_bool
      \bool_set_true:N \l__codedoc_macro_EXP_bool
      \bool_set_false:N \l__codedoc_macro_rEXP_bool
    } ,
  noTF .value_forbidden:n = true ,
  noTF .code:n =
    {
      \bool_set_true:N \l__codedoc_macro_noTF_bool
      \bool_set_true:N \l__codedoc_macro_TF_bool
    } ,
  added .code:n = { \__codedoc_date_set_past:Nn \l__codedoc_date_added_tl {#1} },
  updated .code:n = { \__codedoc_date_set_past:Nn \l__codedoc_date_updated_tl {#1} } ,
  deprecated .bool_set:N = \l__codedoc_macro_deprecated_bool ,
  no-user-doc .bool_set:N = \l__codedoc_macro_nodoc_bool ,
  tested .code:n = { } ,
  label .code:n =
    {
      \clist_set:Nn \l__codedoc_function_label_clist {#1}
      \bool_set_true:N \l__codedoc_no_label_bool
    } ,
  verb .value_forbidden:n = true ,
  verb .bool_set:N = \l__codedoc_names_verb_bool ,
  module .tl_set:N = \l__codedoc_override_module_tl ,
}


\documentclass{l3doc}
\usepackage{lipsum} % to generate some text



\begin{document}
\lipsum


\begin{function}{\macro}
this is a macro.
\end{function}

\begin{environment}{env}
this is an env.
\end{environment}

\begin{function}{opt}
\begin{syntax}
    opt = \meta{a}
\end{syntax}
this is an opt.
\end{function}

\PrintIndex
\end{document}



% ==> transparent bug
\InputIfFileExists{zlatex-cfg.tex}{}{}
\documentclass{../code/zlatex}
% \usepackage{tikz}
% \usepackage{transparent}

\begin{document}
% \zlatexPageMask[anchor=tr, position={(\zpw, \zph)}]{
%   \begin{tikzpicture}
%     \node[opacity=0.5, inner sep=0pt] at (current page.center) 
%       {\includegraphics{Epmmy.jpg}};
%   \end{tikzpicture}
% }

% \zlatexPageMask{\transparent{.1}\includegraphics{Epmmy.jpg}}
% Hello world

\begin{theorem}[AAA]
  This is a theorem AAA.
\end{theorem}

\zlatexThmStyleNew{
  aaa={begin=, end=, option=\zlatexThmTitleSwitch},
}
% \zlatexThmTitleSwitch
\zlatexThmStyle{aaa}

\begin{theorem}[BBB]
  This is a theorem BBB.
\end{theorem}

\end{document}


% \DocumentMetadata{}
% \documentclass{article}
\documentclass{l3doc}

\usepackage{eso-pic,graphicx,transparent,lipsum}

\begin{document}
\lipsum[1-3]

\AddToShipoutPictureFG*{%
  \AtPageLowerLeft{%
    \transparent{.6}\includegraphics[width=\paperwidth,height=\paperheight]{back.jpeg}%
  }%
}

\lipsum[4-6]

\end{document}




% ==> cleverref bug: fixed after updating
\documentclass{book}
\usepackage[UTF8, heading]{ctex}
\usepackage{cleveref}
% \InputIfFileExists{zlatex-cfg.tex}{}{}
% \documentclass[class=book, lang=cn]{../code/zlatex}

\begin{document}
\section{A}
Hello world\label{aaa}

Hello \cref{aaa}
\end{document}



% ==> l3doc cover
\InputIfFileExists{zlatex-cfg.tex}{}{}
\documentclass[lang=cn]{../code/zlatex}


\title{z\LaTeX{} 用户手册}
\author{Eureka}
\date{\today}
\begin{document}
\maketitle
\end{document}




% ==> custom thm style
\InputIfFileExists{zlatex-cfg.tex}{}{}
\documentclass{../code/zlatex}
\usepackage{lipsum}


\begin{document}
\ExplSyntaxOn
\zlatexThmStyleNew {
  test = {
    begin  = \noindent\hspace*{1em}{\zlatexThmTitle}, 
    end    = , 
    option = \__zlatex_thm_title_inline:n {F}
  }
}
\ExplSyntaxOff
\zlatexThmStyle{test}

\section{First}
This is a simple theorem style in z\LaTeX{}:
\begin{theorem}[AAA]
  This is a simple definition. \lipsum[1][1-2]
  \begin{align}
    a^2 + b^2 = c^2
  \end{align}
  \lipsum[1][6-8]
\end{theorem}

\end{document}



\InputIfFileExists{zlatex-cfg.tex}{}{}
\documentclass[
  fancy, 
  % hyper,
  mathSpec={envStyle=elegant}, 
  % layout={slide, aspect=16|9}
]{../code/zlatex}
\usepackage{lipsum}
% \zlatexThmStyle{lapsis}
% \zlatexSetup{mathSpec={envStyle=shadow}}
\zlatexloadlibrary{theme}
\zlatexThmCreate{theorem}{Ztheorem=NEW-THM|{HTML}{8e22ad}, Zproposition} 
% \zlatexThmTitleFormat{\bfseries\zlatexThmName\ \zlatexThmNumber\zlatexThmNote{(}{)}}
% \zlatexThmBefore{\quad}

\def\boomen{As any dedicated reader can clearly see, the Ideal of practical
reason is a representation of, as far as I know, the things in themselves; 
\begin{align}
\underset{}{\mathbf{v} \bigotimes \mathbf{w}} 
  & = \underset{}{\mathbf{v} \otimes \mathbf{w}}
      = \sum_{i=1}^3\sum_{j=1}^3a_{ij}u^iv^j \\[-.75em]
  & = \sum_{i=1}^3\left(a_{i1}u^iv^1+a_{i2}u^iv^2+a_{i3}u^iv^3\right) 
  \end{align}  
As any dedicated reader can clearly see, the Ideal of practical
reason is a representation of, as far as I know, the things in themselves;%
}
% \zlatexThmLang{fr}
\zlatexThmCnt{share, parent=section}
\zlatexThmTocPrefix{\ding{118}\;}
\zlatexColorSetup{
  remark=red!50,
  axiom = gray!40,
  definition  = {RGB}{219,48,122},
  theorem = teal!20
} % OK
\ExplSyntaxOn
\newcommand{\Example}{
\clist_map_inline:nn 
  {axiom, definition, theorem, lemma, corollary, proposition, remark}{
    \lipsum[1][1-2]
    \begin{##1}[\MakeUppercase{##1}]
      \boomen
      % folowing contents for 'lapsis' style
      % \tcblower
      % \boomen
    \end{##1}
    \lipsum[2][1-3]
  }
}
\newcommand{\ExampleII}{
\clist_map_inline:nn 
  {plain, leftbar, background, fancy, shadow, paris, elegant, obsidian}{
    \zlatexThmStyle{##1}
    % \Example  % NOTE: --> lapsis faild for slide note
    % \lipsum[1][1-2]
    \begin{remark}[STYLE:##1]
      \boomen
    \end{remark}
    \begin{proposition}[STYLE:##1]
      \boomen
    \end{proposition}
    \lipsum[2][1-2]
  }
}
\ExplSyntaxOff


\title{DEBUG}
\author{Eureka}
\date{\today}
\begin{document}
% \section{List of Theorems}
% \zlatexThmToc

\section{Example Env}
% \Example

\ExampleII
\end{document}





% new thm them 'lapsis'
\InputIfFileExists{zlatex-cfg.tex}{}{}
\documentclass{../code/zlatex}
\zlatexloadlibrary{theme}
\zlatexThmStyle{lapsis}
\usepackage{lipsum}
\def\AAA{\lipsum[1][1-3]
\begin{align}
  \int u\;\mathrm{d}v = uv - \int v \;\mathrm{d}u
\end{align}
\lipsum[1][1-3]}
\def\BBB{\lipsum[1][3-5]
\begin{align}
  {\frac {\mathrm{d}}{\mathrm{d}x}}{\Big (}u(x)v(x){\Big )}=v(x){\frac {\mathrm{d}}{\mathrm{d}x}}
    \left(u(x)\right)+u(x){\frac {\mathrm{d}}{\mathrm{d}x}}\left(v(x)\right)
\end{align}
\lipsum[1][1-3] }

\begin{document}
\section{Lapsis thm Theme}
\lipsum[1-2]
\begin{theorem}[TEST THM]
  \AAA 
  \tcblower
  \BBB
\end{theorem}

\lipsum[3]
\begin{definition}
\AAA
\end{definition}
\lipsum[4-7]

\begin{proposition}[PROP NEW]
  \tcblower
  \BBB
\end{proposition}
\end{document}



% ==> thm backgroud to paperwidth
\InputIfFileExists{zlatex-cfg.tex}{}{}
\documentclass{../code/zlatex}
\usepackage[many]{tcolorbox}
\tcbuselibrary{skins}
\usepackage{lipsum}
\usepackage{xcolor}
\usepackage{tikz}
\usepackage{adjustbox}
\DeclareMathSymbol{\blacktriangleleft}{\mathrel}{AMSa}{"4A}
\usetikzlibrary{fadings, calc}
\parindent0pt
\definecolor{shadecolor}{HTML}{bfbfbf}
\newenvironment{back}{
  \def\FrameCommand{\colorbox{shadecolor}}
  \MakeFramed{\advance\hsize-\width \FrameRestore}
}{\endMakeFramed }
\newcommand{\llapnote}[1]{
  \mbox{}\llap{%
  \adjustbox{set height=0pt, set depth=0pt}{%
    \parbox[t]{2.85cm}{\raggedleft #1}}\hspace*{.75em}}%
}


\begin{document}
% \begin{tcolorbox}[width=\textwidth, arc=0pt]
\lipsum[1]
% \end{tcolorbox}

\zlatexPageMask[
  anchor=tl,
  % layer=foreground,
  position={(0pt, .669\zph)}
]{\tikz\fill[color=gray!50, path fading=east] (0, 0) rectangle +(\zpw, 20em);}
% \begin{tcolorbox}[
%   standard jigsaw,
%   width=\paperwidth, 
%   arc=0pt, opacityback=0,
% ]
%   \lipsum[1]
% \end{tcolorbox}

\begin{back}%
% \adjustbox{set height=0pt, set depth=0pt}{%
%   \tikz[remember picture, overlay, baseline]%baseline=(current bounding box.north), 
%     \fill[color=red!50, path fading=east] (-3.5cm, 0pt) rectangle +(\zpw, -20em);
% }
\llapnote{\sffamily\textbf{Example} 1.1\\Integrate by Part}%
% \llapnote{\sffamily\textbf{Example}\\ 1.1|Integrate by Part\\HELLO WORLD}%
\lipsum[1][1-3]
\begin{align}
  \int u\;\mathrm{d}v = uv - \int v \;\mathrm{d}u
\end{align}
\lipsum[1][1-3]
\end{back}
\lipsum[1][3-5]
\begin{align}
  {\frac {\mathrm{d}}{\mathrm{d}x}}{\Big (}u(x)v(x){\Big )}=v(x){\frac {\mathrm{d}}{\mathrm{d}x}}
    \left(u(x)\right)+u(x){\frac {\mathrm{d}}{\mathrm{d}x}}\left(v(x)\right)
\end{align}
\lipsum[1][1-3]\hfill$\mathcolor{gray}{\blacktriangleleft}$


\newpage
\section{Tcolor box}
Hello world 
\vspace*{6em}

\begin{tcolorbox}[
  enhanced,  breakable,
  top=1.5pt, bottom=1.5pt,
  left=2pt,  leftlower=-3pt,
  right=3pt, arc=0pt, frame hidden,
  bicolor, colback=shadecolor,
  opacitybacklower=0,
  frame code app={
    \draw[color=gray, thick] 
      (frame.north west)++(-5cm, -1pt)--($(frame.north east)+(5cm, 0pt)$);
    \draw[color=gray, thick] 
      (frame.south west)++(-5cm, -1pt)--($(frame.south east)+(5cm, 0pt)$);
    \fill[color=gray!50, path fading=east] 
      (frame.north west)++(-5cm, -1pt) rectangle ($(frame.south east)+(5cm, 0pt)$);
  },
]\llapnote{\sffamily\textbf{Example} 1.1\\Integrate by Part}%
\lipsum[1][1-3]
\begin{align}
  \int u\;\mathrm{d}v = uv - \int v \;\mathrm{d}u
\end{align}
\lipsum[1][1-3]
\tcblower
\lipsum[1][3-5]
\begin{align}
  {\frac {\mathrm{d}}{\mathrm{d}x}}{\Big (}u(x)v(x){\Big )}=v(x){\frac {\mathrm{d}}{\mathrm{d}x}}
    \left(u(x)\right)+u(x){\frac {\mathrm{d}}{\mathrm{d}x}}\left(v(x)\right)
\end{align}
\lipsum[1][1-3]\hfill$\mathcolor{gray}{\blacktriangleleft}$
\end{tcolorbox}


\newpage
\section{NEWSEC}
\subsection{parbox}
Hello world\parbox[t]{6em}{aaaaaa bbbbb cccc}


% Hello world \tikz[anchor=north, baseline=(A.base)]\fill[color=red!50, path fading=east] (0, 0)++(-3cm, 0pt) rectangle +(5em, 8em);

\subsection{tikz object}
before text 
\tikz[baseline={([yshift={-\ht\strutbox}]a.north)},outer sep=0pt,inner sep=0pt] \node[align=center] (a) {\strut line1\\line2\\line3\\line4};
text between
\tikz[baseline=(a.south),outer sep=0pt,inner sep=0pt] \node[align=center] (a) {\strut line1\\line2\\line3\\line4};
text after


\subsection{tcb upper/lower}
\newtcolorbox{BBB}{breakable,
enhanced,  breakable,
  bicolor,
  colback=red!10!white,
  colbacklower=blue!5!white,
  title={Example}}
\begin{BBB}
  \lipsum[1]
  \tcblower
  \lipsum[2]
\end{BBB}
\end{document}




% ==> test obsidian thm style 
\InputIfFileExists{zlatex-cfg.tex}{}{}
\documentclass[layout={slide, aspect=16|9}]{../code/zlatex}
\usepackage[many]{tcolorbox}
\zlatexloadlibrary{theme}
\zlatexThmStyle{obsidian}
\usepackage{lipsum}

\title{test OBSIDIAN THM}
\author{Eureka}
\date{\today}
\begin{document}
\maketitle
\section{FISRT}
\lipsum[1] 
\begin{theorem}[THM]
  \lipsum[1][1-5] 
  \[
    \sum_{i=1}^{+\infty}{\int_{0}^{i}-\frac{1}{t}\mathrm{d}t} = \frac{\pi^2}{6}
  \]
\end{theorem}

\lipsum[2] 
\begin{proposition}[PROP]
  \lipsum[1][1-5] 
  \[
    \sum_{i=1}^{+\infty}{\int_{0}^{i}-\frac{1}{t}\mathrm{d}t} = \frac{\pi^2}{6}
  \]
\end{proposition}

\lipsum[3] 
\begin{definition}
  \lipsum[1][1-5] 
  \[
    \sum_{i=1}^{+\infty}{\int_{0}^{i}-\frac{1}{t}\mathrm{d}t} = \frac{\pi^2}{6}
  \]
\end{definition}
\end{document}


% ==> '+=' keys feature
\InputIfFileExists{zlatex-cfg.tex}{}{}
\documentclass{../code/zlatex}
\parindent=0pt

\begin{document}
\ExplSyntaxOn
\keys_define:ne {test}{
  % simple key
  keyA    .tl_set:N = \exp_not:c {l__test_keyA_tl},
  keyA+   .code:n   = \tl_put_right:Nn \exp_not:c {l__test_keyA_tl} {#1},
  keyA~+  .code:n   = \tl_put_right:Nn \exp_not:c {l__test_keyA_tl} {#1},
  % zlatex internal
  keyB    .meta:nn  = { text / keyB }{#1},
  \__zlatex_plus_key_aux:nnn {l__test_keyB_tl}{keyB}{keyB-x},
  \__zlatex_plus_key_aux:nnn {l__test_keyB_tl}{keyB}{keyB-y},
}
\NewDocumentCommand{\testPlusKeyA}{m}{
  \keys_set:nn {test}{#1}
  \tl_use:N \l__test_keyA_tl\par
}
\NewDocumentCommand{\testPlusKeyB}{m}{
  \keys_set:nn {test}{#1}
  \tl_use:N \l__test_keyB_tl\par
}

\ExplSyntaxOff
\section{Simple}
\testPlusKeyA{keyA = AA}
\testPlusKeyA{keyA+ = aa} 
\testPlusKeyA{keyA += xx} 
\testPlusKeyA{keyA=zz}

\section{Internal}
\testPlusKeyB{keyB/keyB-x = BB}
\testPlusKeyB{keyB/keyB-x+ = bb}
\testPlusKeyB{keyB/keyB-x += xx}
\testPlusKeyB{keyB/keyB-x = zz}

\dotfill\par
\testPlusKeyB{keyB/keyB-y = BBYY}
\testPlusKeyB{keyB/keyB-y+ = bbyy}
\end{document}





% ==> mathalias library/pagemask test
\InputIfFileExists{zlatex-cfg.tex}{}{}
\documentclass{../code/zlatex}
\zlatexloadlibrary{mathalias}
\zlatexThmStyle{fancy}
\usepackage{tikz}
\usetikzlibrary{fadings}
\usepackage{lipsum}


\begin{document}
\zlatexPageMask[
  anchor=bl,
  position={(0pt, -14em+\zph)}
]{\includegraphics[width=\zpw]{back.jpeg}}
\zlatexPageMask[
  % label=FIRST, % EXPECT: Class zlatex Warning: Only star version of \zlatexPageMask is label-allowed.
  anchor=bl,
  position={(0pt, -14em+\zph)}
]{%
  \tikz\fill[color=gray, path fading=north] (0, 0) rectangle +(\zpw, 5em);
}
\thispagestyle{empty}
{
  \def\thesection{\textcolor{white}{\arabic{section}}}
  \section{\textcolor{white}{Symbols}}
}\vspace*{3em}
\subsection{simple}
\lipsum[1]
\begin{align}
  & \la \\
  & \rra \\
  & \Nra \\
  & \dda \\
  & \xra*[Up](down)\\
  & \xla(\beta) \\
  & \xxla[2](1)
\end{align}

\subsection{Complex}
Some Examples: $H^1(\Omega)\hra L^p(\Omega)$, And further more: 
\begin{align}
  H^1(\Omega)\hra[E:t-1 = \gamma] L^p(\Omega)
\end{align}

Some symbols: \CC, \RR, \ZZ, \NN

\begin{align}
  & \zab \alpha + \beta \text{ V.S. } \alpha + \beta \\
  & \zab(\frac12) \text{ V.S. } (\frac12) \\
  & \zab[\frac12] \text{ V.S. } [\frac12] \\
  & \zab\{\frac12\} \text{ V.S. } \{\frac12\}
\end{align}

\newpage
\zlatexPageMask*[anchor=tc,label=TEST,layer=foreground]{%
  % \tikz\fill[color=blue, path fading=north] (0, 0) rectangle +(\zpw, 5em);
  \tikz\fill[semitransparent, teal] (-10em,-10em) rectangle +(10em,10em); 
}
\section{Text}
\lipsum[1]

\begin{definition}[Test THM]
  A smooth vector field \(\mathbf{v}\) is defined in a domain \(M\) if to each point \(x\) there is 
  assigned a vector \(\mathbf{v}\left( x\right)  \in  {T}_{x}M\) attached at that point and depending 
  smoothly on the point \(x\) (if a system of \(m\) coordinates is chosen, the field is defined by 
  its \(m\) components, which are smooth functions of \(m\) variables). The vector \(\mathbf{v}\left( x\right)\) 
  is called the value of the field \(\mathbf{v}\) at the point \(x\) .
\end{definition}

\lipsum[1]

\begin{definition}[Continuous Mapping]
  A smooth vector field \(\mathbf{v}\) is defined in a domain \(M\) if to each point \(x\) there is 
  assigned a vector \(\mathbf{v}\left( x\right)  \in  {T}_{x}M\) attached at that point and depending 
  smoothly on the point \(x\) (if a system of \(m\) coordinates is chosen, the field is defined by 
  its \(m\) components, which are smooth functions of \(m\) variables). The vector \(\mathbf{v}\left( x\right)\) 
  is called the value of the field \(\mathbf{v}\) at the point \(x\) .
\end{definition}

\lipsum[3]

\begin{lemma}[Test THM]
  A smooth vector field \(\mathbf{v}\) is defined in a domain \(M\) if to each point \(x\) there is 
  assigned a vector \(\mathbf{v}\left( x\right)  \in  {T}_{x}M\) attached at that point and depending 
  smoothly on the point \(x\) (if a system of \(m\) coordinates is chosen, the field is defined by 
  its \(m\) components, which are smooth functions of \(m\) variables). The vector \(\mathbf{v}\left( x\right)\) 
  is called the value of the field \(\mathbf{v}\) at the point \(x\) .
\end{lemma}

\lipsum[4]
\zlatexPageMaskRemove{foreground}{TEST}
\end{document}


% ==> Fixed Bugs: 
% 1. THM note expansion bug; 
% 2. hyper init faild; 
% 3. toc line stretch 
%    --> fixed bug: line-stretch failed when no \tableofcontents before
\InputIfFileExists{zlatex-cfg.tex}{}{}
\documentclass[hyper]{../code/zlatex}
\zlatexloadlibrary{mathalias}
\zlatexColorSetup{link=teal}
\zlatexSetup{
  toc={
    stretch=2, 
    title-vspace=-5em, 
    title=\centerline{\bfseries\Large \textsc{Contents}}
  }
}
\usepackage{lipsum}
\zlatexThmBefore{\par\noindent\dotfill\par}
\zlatexThmTocPrefix{\P\;}
% \zlatexThmTocSymbolClear
\zlatexThmTocSymbol{
  lemma=$\star$\;,
  definition=$\diamond$\;,
  theorem=$\clubsuit$\;,
}
\ExplSyntaxOn
\newcommand{\addthmsec}[4]{
  \__zlatex_thm_toc_add:nnnn {#1}{#2}{#3}{#4}
}
\ExplSyntaxOff


\begin{document}
\thispagestyle{empty}
\tableofcontents
\vskip5em

\lipsum[1]
\section{List of Theorems}
Lorem ipsum dolor sit amet, consectetuer adipiscing elit. Ut purus elit, vestibulum
ut, placerat ac, adipiscing vitae, felis.

\zlatexThmToc[3]

\section{Accent}
\addthmsec{section}{BLANK}{\thepage}{section.1}
\lipsum[2]
\begin{lemma}[N\"ormal Text]
  This is a simple lemma.
  \begin{align}
    a^2 + b^2 = c^2
  \end{align}
\end{lemma}

\lipsum[2]
\begin{proposition}[Normal Text]
  This is a simple lemma.
  \begin{align}
    a^2 + b^2 = c^2
  \end{align}
\end{proposition}


\section{Macros/Equation}
\zlatexThmTocAdd[name=NEW-SEC]{section}
\begin{theorem}[$\mathbb{Z}[\sqrt{2}]$]%
  This is a simple theorem.
  \begin{align}
    a^2 + b^2 = c^2
  \end{align}
\end{theorem}


\begin{definition}
  This is a simple definition.
  \begin{align}
    a^2 + b^2 = c^2
  \end{align}
\end{definition}

\ExplSyntaxOn
% \cs_show:N \sqrt

\cs_set:Npn \__test_sqrt:n #1 
  { #1--$Z[\sqrt {2}]$ }

\__test_sqrt:n {AA}

\the\baselineskip
\ExplSyntaxOff


\section{AAA}
\subsection{AAA-1}
\subsection{AAA-2}
\subsection{AAA-3}

\section{BBB}
\subsection{BBB-1}
\subsection{BBB-2}
\subsection{BBB-3}
\end{document}


% ==> thmnote update
\InputIfFileExists{zlatex-cfg.tex}{}{}
\documentclass{../code/zlatex}


\begin{document}
\section{First}
This is a simple definition style in z\LaTeX{}:

\begin{definition}
  This is a simple definition. 
  \begin{align}
    a^2 + b^2 = c^2
  \end{align}
\end{definition}

A simple theorem environment follows the above definition as below:
\begin{theorem}[THM-NOTE]
  This is a simple theorem.
  \begin{align}
    a^2 + b^2 = c^2
  \end{align}
\end{theorem}


\section{SET to HEIGHT}
\ExplSyntaxOn
\ztool_resize_to_ht:nn {2.5pt}{XXX}
\ExplSyntaxOff


\section{THM LIST}
Test List of Theorems:

\zlatexThmToc
\end{document}





% ==> slide mode test
\InputIfFileExists{zlatex-cfg.tex}{}{}
\documentclass[
  hyper,
  % class=book,
  % lang=cn,
  layout={slide, aspect=16|9, theme=AnnArborSpruce},
]{../code/zlatex}
\zslideThemeUse[
  UR={text=\zslideIfPageTF{=1}{}{\zslideDefaultUR:\ \zslideNavigateBall}},
]{AnnArborSpruce}
\makeatletter
\def\seeTargets{%
  \par\noindent\dotfill\\
  \string\@currentHref\ = \meaning\@currentHref\par
  \string\@currentHpage\ = \meaning\@currentHpage\par
  \string\Hy@currentbookmarklevel\ = \meaning\Hy@currentbookmarklevel\par\vskip3em
}
\makeatother
\zlatexPageMask*[layer=foreground]{AAA}
\zlatexPageMask[layer=foreground]{BBB}
% \zslidelogo[position={(40pt, 10pt)}, exclude={0, 2}]{Epmmy.jpg}
\zslidelogo[width=3em, exclude={0, 1}]{Epmmy.jpg}


\title{Debug}
\author{Eureka}
\date{\today}
\begin{document}
\maketitle
\tableofcontents

\section{First Section}
\subsection{Test}
BULLET TEST: \textbullet

\ExplSyntaxOn
m\ztool_resize_to_wd:nn {1em}{mmmm}\par 

\rule{1em}{3pt}
\ExplSyntaxOff

\zslideIfPageTF{<2}{PAGE-II}{NOT PAGE-II}

\newpage
\subsection[Short First Subsection Name]{First Subsection Name}
  \begin{itemize}
    \item
    An item
    \item
    Another item
    \begin{itemize}
      \item Something
      \begin{itemize}
        \item Something else
      \end{itemize}
    \end{itemize}
    \begin{enumerate}
      \item Thing A
      \item Thing B
      \item Thing C
    \end{enumerate}
  \end{itemize}
\end{document}




% ===> Re-write titlesec package
\InputIfFileExists{zlatex-cfg.tex}{}{}
\documentclass[
  % hyper,
  class=book,
]{../code/zlatex}
\makeatletter
% --> LaTeX 2e interface
\newcounter{sector}
\newcommand{\sectormark}[1]{\markright{#1}}
\newcommand{\sector}{\@startsection{sector}{1}{0pt}%
  {-3.5ex plus -1ex minus -.2ex}
  {2.3ex plus.2ex}
  {\large\sffamily}% start normal sec syntax: "*[]{}"
}
\makeatother

% --> zlatex interface
% \zlatexTitleNew{page}{part}



% \zlatexTitleStyle{section}
% {
%   shape=display,
%   % align=left,
%   label=X-\thesection-Y.,
%   sep=10pt,
%   format=\color{blue},
%   labelformat=\bfseries,
%   titleformat=\sffamily\color{red},
%   titlecmd=\testcmd,
%   before=\raggedleft\titlerule[1pt]\vskip2em|,
%   after=\vskip3em,
% }
% \ExplSyntaxOn
% \group_begin:
% \__zlatex_title_format_copy:nnnnnnn 
%   {\section}{hang}{\sffamily}{\thesection~ XX}{0pt}{}{}
% \group_end:
% \ExplSyntaxOff


\begin{document}
% \part{PART}
\chapter{CHAPTER}
\section{AAA}
The prior examples redefined existing sectional unit title 
commands. This defines a new one, illustrating the needed 
counter and macros to display that counter.


\sector{BBB}
The prior examples redefined existing sectional unit title 
commands. This defines a new one, illustrating the needed 
counter and macros to display that counter.



\begin{align}
  \substack{\smash{A^T}\\\scriptscriptstyle n\text{by}m}
  \substack{\smash{C}\\\scriptscriptstyle m\text{by}m}
  \substack{\smash{A}\\\scriptscriptstyle m\text{by}n}
  = 1
\end{align}
\end{document}





% ==> learn titlesec
\documentclass{article}
\usepackage[hmargin=2in]{geometry}
\usepackage[raggedleft, sf]{titlesec}
\titlelabel{\thetitle.\;}
\usepackage{xcolor}

% \newcounter{part}
\titleclass{\part}{straight}
\titleformat{\part}{\bfseries}{\thepart.}{10pt}{}

% 1. \part is usually implemented in a non-standard way, it 
%    remains untouched by the simple settings and should be 
%    changed with the help of the “Advanced Interface.”


\begin{document}
\part{First Part}
\section{FIRST}
The prior examples redefined existing sectional unit title 
commands. This defines a new one, illustrating the needed 
counter and macros to display that counter.


\section{SECOND}
The prior examples {redefined existing} sectional unit title.
commands. This defines a new one, illustrating the needed 
counter and macros to display that counter. commands This 
defines a new one, illustrating the needed counter and macros 
to display that counter illustrating the needed defines a new one, 
illustrating the needed counter and macros to display that counter.

\ExplSyntaxOn
\cs_new:Nn \_test:nn {
  #1 == #2
}
\_test:nn {AA}{BB}
\ExplSyntaxOff

\end{document}




% ===> page info
\InputIfFileExists{zlatex-cfg.tex}{}{}
\documentclass[
  hyper
]{../code/zlatex}
\usepackage[many]{tcolorbox}
\makeatletter
\zlatexPageMask[
  anchor=tr, 
  position={(\zpw, \zph)}
]{\color{orange}\rule{\zpw}{10pt}}
\colorlet{colexam}{red!75!black}
\zlatexThmCnt{share, parent=subsection}

\begin{document}
HELLO WORLD \rule{10pt}{10pt}

\section{TEXT}
\zlatexPageMask{AAA}
\zlatexPageMask{\includegraphics{Epmmy.jpg}}
\zlatexPageMask[layer=foreground, position={(10em, 10em)}]{\parbox{3em}{BBB-1\\BBB-2}}
\zlatexPageMask*[position={(.25\zpw, .25\zph)}]{\ifnum\thepage=2 CCC2\else CCC1\fi}
\zlatexPageMask*[position={(0pt, 0pt)}, anchor=lb]{XXX-\thepage}
\zlatexPageMask[position={(\zpw, \zph)}, anchor=tr]{YYY}

\newpage
\zlatexPageMask[
  anchor=br, 
  position={(\zpw, 0pt)}
]{\color{purple}\rule{15pt}{\zph}}


\section{GRAPH}
\subsection{AAA}
% ==> style definitions
\tcbset{
  styleA/.style={
    arc=0mm, 
    bottomtitle=0.5mm,
    boxrule=0mm,
    colbacktitle=black!10!white, 
    coltitle=black, 
    fonttitle=\bfseries, 
    left=2.5mm,
    leftrule=1mm,
    right=3.5mm,
    title={\zlatexThmTitle},
    toptitle=0.75mm, 
  },
  styleB/.style={
    empty,
    frame engine=path,
    colframe=yellow!10,
    sharp corners,
    title={\zlatexThmTitle},
    attach boxed title to top left={yshift*=-\tcboxedtitleheight},
    boxed title style={size=minimal, top=4pt, left=4pt},
    coltitle=colexam,fonttitle=\large\bfseries\sffamily,
    drop fuzzy shadow,
    coltitle=black,
    borderline west={3pt}{-3pt}{teal!50},
    attach boxed title to top left={xshift=-3mm, yshift*=-\tcboxedtitleheight/2},
    boxed title style={right=3pt, bottom=3pt, overlay={
      \draw[draw=teal!70, fill=teal!70, line join=round]
        (frame.south west) -- (frame.north west) -- (frame.north east) --
        (frame.south east) -- ++(-2pt, 0) -- ++(-2pt, -4pt) --
        ++(-2pt, 4pt) -- cycle;
    }},
    overlay unbroken={
      \scoped \shade[left color=teal!10!black, right color=teal]
        ([yshift=-0.2pt]title.south west) -- ([xshift=-1.5pt, yshift=-0.2pt]%
        title.south-|frame.west)-- ++(0, -6pt) -- cycle;
    },
  },
  styleC/.style={
    enhanced, 
    breakable,
    colback=white,
    colframe=blue!30!black,
    attach boxed title to top left={yshift*=-\tcboxedtitleheight}, 
    title={\zlatexThmTitle},
    boxed title size=title,
    boxed title style={%
        sharp corners, 
        rounded corners=northwest, 
        colback=tcbcolframe, 
        boxrule=0pt,
    },
    underlay boxed title={%
        \path[fill=tcbcolframe] (title.south west)--(title.south east) 
            to[out=0, in=180] ([xshift=5mm]title.east)--
            (title.center-|frame.east)
            [rounded corners=\kvtcb@arc] |- 
            (frame.north) -| cycle; 
    },
  }
}
\zlatexThmStyleNew{
  % ---->
  styleA={
    begin={\begin{tcolorbox}[styleA]},
    end={\end{tcolorbox}}
  },
  % ---->
  styleB={
    begin={\begin{tcolorbox}[styleB]},
    end={\end{tcolorbox}}
  },
  % ---->
  styleC={
    begin={\begin{tcolorbox}[styleC]},
    end={\end{tcolorbox}}
  }
}
% ===> test text
\def\bench{Mathematics is a field of study that discovers and organizes methods, 
theories and theorems that are developed and proved for the needs of empirical 
sciences and mathematics itself. There are many areas of mathematics, 
\begin{align}
  \mathrm{d}^2x + \mathrm{d}^y2 = \mathrm{d}^2z
\end{align}
which include number theory (the study of numbers), algebra (the study of formulas and related
structures), geometry (the study of shapes and spaces that contain them), analysis 
(the study of continuous changes), and set theory (presently used as a 
foundation for all mathematics).}

% \zlatexThmStyle{styleA}
\zlatexThmTitleFormat*{%
  \bfseries\color{orange}%
  \zlatexThmNote\ \zlatexThmNumber\ \zlatexThmName
}
\begin{definition}[DEF]
\bench
\end{definition}


% \zlatexThmStyle{styleB}
\vskip8em
\begin{theorem}[THM]
\bench
\end{theorem}


\newpage
% \zlatexThmStyle{styleC}
\begin{axiom}[AXIOM]
\bench
\end{axiom}


\newpage
\zlatexThmToc
\end{document}






% ===> keytheorems ref
\InputIfFileExists{zlatex-cfg.tex}{}{}
\documentclass{article}
\usepackage{keytheorems}
\newkeytheorem{definition}
\usepackage[colorlinks]{hyperref}

\begin{document}
\stepcounter{page}
\begin{definition}[DEF]\label{def:a}
Hello world, this is a definition. $a^2 + b^2 = \mathrm{d}^2z$
\end{definition}


\listofkeytheorems


\newpage
\listofkeytheorems

\@writefile{thlist}{\KeyThmsSavedTheorem{definition}{1}{definition.1}{2}{}{note={DEF}}{}}
\KeyThmsSavedTheorem {definition}{1}{}{1}{}{note={DEF}}{}

\cs_new_protected:Npn \__keythms_thm_addcontentsdata:nnnn #1#2#3#4
  { % #1 = theorem name, #2 = stored counters, #3 = keys, #4 = body
    \keythms_listof_chaptervspacehack:
    \iow_shipout:Ne \@auxout
      {
        \token_to_str:N \@writefile { thlist }
          {
            \token_to_str:N \KeyThmsSavedTheorem{ #1 }
              { \@currentlabel }
              { \@currentHref }
              { \thepage }
              { #2 }
              { \exp_not:n { #3 } } % do we want any expansion here, perhaps
              { \exp_not:n { #4 } } % with \text_expand:n ?
          }
      }
  }
\end{document}




% ==> add list of theorem interface
\InputIfFileExists{zlatex-cfg.tex}{}{}
\documentclass[hyper]{../code/zlatex}
\makeatletter
\usepackage[many]{tcolorbox}

\zlatexThmTocLevel{subsection}
\zlatexThmStyle{shadow}
\begin{document}
\zlatexThmToc

\vskip6em
\subsection{AAA}
\begin{definition}[DEF]\label{def:aaa}
  Hello world, this is a definition. $a^2 + b^2 = \mathrm{d}^2z$
\end{definition}


\newpage
\subsection{BBB}
\begin{theorem}[THM]
  Hello world, this is a theorem. $a^2 + b^2 = \mathrm{d}^2z$
\end{theorem}


\newpage
\zlatexThmToc[vspace=15pt, title=THM LIST]
\end{document}


\InputIfFileExists{zlatex-cfg.tex}{}{}
\documentclass[hyper]{../code/zlatex}
\parindent=0pt \makeatletter
\ExplSyntaxOn
% \NewHook{zlatex/thm/thlist}
\def\@starttoc#1{%
  \begingroup
  \@input{\jobname.#1}
  % \@fileswfalse
  % \UseHook{zlatex/thm/thlist}
  % \legacy_if_set_false:n { @filesw }
  \typeout{------->\checksave}
  \if@filesw
    \typeout{----->enter~stream}
    \expandafter\newwrite\csname tf@#1\endcsname
    \immediate\openout \csname tf@#1\endcsname \jobname.#1\relax
  \fi
  \@nobreakfalse
  \endgroup
}


% optional commands
\def\seeTargets{%
  \par\noindent\dotfill\\
  \string\@currentHref\ = \meaning\@currentHref\par
  \string\@currentHpage\ = \meaning\@currentHpage\par
  \string\Hy@currentbookmarklevel\ = \meaning\Hy@currentbookmarklevel\par\vskip3em
}
\def\checksave{
  \if@filesw SAVED \else NOT SAVED \fi
}
\ExplSyntaxOff

\begin{document}
\checksave
\vskip3em

\@starttoc{thlist}
% \@@input debug.thlist
\newpage

\section{FIRST}
\seeTargets

\begin{definition}[DEF]\label{def:aa}
Hello world, this is a definition. $a^2 + b^2 = \mathrm{d}^2z$
\end{definition}
\addcontentsline{thlist}{subsection}{DEF-ITEM}
\seeTargets

\section{THM-LIST}
\meaning\@@input

% \checksave
\begingroup
\@fileswfalse
\@starttoc{thlist}
% \input{\jobname.thlist}
\@@input debug.thlist
\endgroup
% \contentsline {subsection}{DEF-ITEM}{2}{definition.1.1}%


\end{document}







% ==> New thm style interface
\InputIfFileExists{zlatex-cfg.tex}{}{}
\documentclass[
  fancy, 
  hyper, 
  % layout={slide, aspect=16|9}
]{../code/zlatex}
\usepackage{pifont}
\zlatexSetup{mathSpec={envStyle=shadow}}
\zlatexThmCreate{theorem}{Ztheorem=NEW-THM|{HTML}{8e22ad}, Zproposition} 


\def\boomen{As any dedicated reader can clearly see, the Ideal of practical
reason is a representation of, as far as I know, the things in themselves; 
\begin{align}
\underset{}{\mathbf{v} \bigotimes \mathbf{w}} 
  & = \underset{}{\mathbf{v} \otimes \mathbf{w}}
      = \sum_{i=1}^3\sum_{j=1}^3a_{ij}u^iv^j \\[-.75em]
  & = \sum_{i=1}^3\left(a_{i1}u^iv^1+a_{i2}u^iv^2+a_{i3}u^iv^3\right) 
  \end{align}  
}
\zlatexThmLang{fr}
\zlatexThmCnt{share, parent=section}
\zlatexThmTocPrefix{\ding{118}\;}
\zlatexColorSetup{
  definition  = {RGB}{219,48,122},
  theorem = teal!20
} % OK

\title{DEBUG}
\author{Eureka}
\date{\today}
\begin{document}
% BUG: stretch not working in thm toc.
% \maketitle
% \zlatexThmToc[stretch=20]
% \zlatexThmToc[
%   title={\textsc{List of Theorems}},
%   stretch=1.5, 
%   title-vspace=-10pt, 
%   after-vspace=20pt
% ]


\section{NEWLY}
\begin{Ztheorem}[ZTHM]
\boomen
\end{Ztheorem}

\zlatexThmColorSetup{
  Zproposition={HTML}{c0eb75},
}
\zlatexThmStyle{leftbar}
\begin{Zproposition}[ZPROP]
\boomen
\end{Zproposition}


\section{DEBUG}
\subsection{AAA}
% \zlatexColorSetup{
%   definition  = {RGB}{219,48,122},
% } % NO effect
\zlatexThmStyle{fancy}
\begin{definition}[DEF]\label{def:1}
\boomen
\end{definition}
\zlatexThmStyleNew{
  aaa={begin=\par\begin{center}, end=\end{center}, option=\zlatexThmTitleSwitch},
  bbb={begin=BBB\begin{center}, end=\end{center}\kaaklfg\ahjksdh\par},
}
% \ExplSyntaxOn
% \cs_meaning:N \g__zlatex_thm_style_aaa_begin_tl
% \cs_meaning:N \g__zlatex_thm_style_aaa_end_tl

% \cs_meaning:N \g__zlatex_thm_style_bbb_begin_tl
% \cs_meaning:N \g__zlatex_thm_style_bbb_end_tl
% \ExplSyntaxOff


\zlatexThmStyle{background}
% \ShowHook{zlatex/thm/titleformat}
\zlatexThmTitleFormat{\bfseries\zlatexThmNumber\ \zlatexThmName}
% \ShowHook{zlatex/thm/titleformat}
\begin{theorem}[THM]\label{thm:2}
  \boomen
\end{theorem}
% \ShowHook{zlatex/thm/titleformat}

\subsection{CCC}
\zlatexThmStyle{elegant}
% \ShowHook{zlatex/thm/titleformat}
\begin{proposition}[PROP]\label{prop:5}
  \boomen
\end{proposition}
% \ShowHook{zlatex/thm/titleformat}

\zlatexThmStyle{aaa}
\section{BBB}
\zlatexThmColorSetup{
  definition={HTML}{22b8cf},
  proof=orange
}
\zlatexThmHook{before=\color{blue}\begin{framed}, after=\end{framed}}
\begin{definition}[DEF-II]\label{def:3}
\boomen
\end{definition}

\zlatexThmStyle{paris}
\begin{definition}[DEF-III]\label{def:4}
  \boomen
\end{definition}


\zlatexThmStyle{leftbar}
\begin{theorem}[REMARK]\label{rem:1}
  As any dedicated reader can clearly see, the Ideal of practical
  reason is a representation of, as far as I know, the things in themselves;
\boomen
As any dedicated reader can clearly see, the Ideal of practical
reason is a representation of, as far as I know, the things in themselves;
\end{theorem}


\section{PROOF-LIKE}
\begin{proof}
Hello this is a simple proof env test.
\end{proof}


\cref{def:1,prop:5} | \Cref{thm:2} | \cref{def:3,rem:1}
\begin{axiom}[AXIOM]
  \boomen
\end{axiom}
\end{document}











% Marker in zslide bug fixed
\InputIfFileExists{zlatex-cfg.tex}{}{}
\documentclass[
  hyper,
  fancy,
  layout={slide, aspect=16|9, theme=AnnArborBeaver}
]{../code/zlatex}
\zlatexSetup{mathSpec={envStyle=fancy}}
\zlatexThmCreate{theorem}{Ztheorem=THM|{HTML}{8e44ad}, Zproposition} 
% \zslideThemeUse[
%   UR={text=\zslideNavigateBall},
% ]{AnnArborBeaver}
\NewMarkClass{test}
\usepackage{lipsum}
% \usepackage{minted}
\def\boomen{As any dedicated reader can clearly see, the Ideal of practical
reason is a representation of, as far as I know, the things in themselves; 
\begin{align}
\underset{}{\mathbf{v} \bigotimes \mathbf{w}} 
  & = \underset{}{\mathbf{v} \otimes \mathbf{w}}
      = \sum_{i=1}^3\sum_{j=1}^3a_{ij}u^iv^j \\[-.75em]
  & = \sum_{i=1}^3\left(a_{i1}u^iv^1+a_{i2}u^iv^2+a_{i3}u^iv^3\right) 
  \end{align}  
}


\ExplSyntaxOn\makeatletter
\makeatother\ExplSyntaxOff

\title{ltMarks Test}
\author{Eureka}
\date{\today}
\begin{document}
\maketitle

% blank page
% \newpage

\section{HELLO}
\InsertMark{test}{AAA-1}
\subsection{WORLD}
\lipsum[1][1-5]
\InsertMark{test}{AAA-2}
\vspace*{8em}

\subsection{SPACE}
\lipsum[1][6-15]

\subsection{ROOM}
\lipsum[1][13-17]
\vspace*{6em}

\subsection{OCEAN}
\lipsum[2][1-6]
\vspace*{10em}

\subsection{HEAVEN}
\lipsum[2][1-16]

% \newpage
% aa


\section{DEBUG}
\subsection{AAA}
\lipsum[1][1-5]
\InsertMark{test}{AAA-2}
\vspace*{8em}

\subsection{BBB}
\lipsum[1][6-15]


\section{FIRST}
\begin{definition}[DEF]
  Hello world: $a^2+b^2=c^2$.
\end{definition}

\zlatexThmStyle{elegant}
\begin{theorem}[theorem Test]
  Hello world: $a^2+b^2=c^2$.
\end{theorem}


\section{SECOND}
% \zlatexColorSetup{
%   % Ztheorem={HTML}{8e44ad},
%   Ztheorem=red 
% }
\begin{Ztheorem}[Ztheorem Test]
  Hello world: $a^2+b^2=c^2$.
\end{Ztheorem}

\begin{Zproposition}
  Hello world: $a^2+b^2=c^2$.
\end{Zproposition}


\newpage
% \zlatexColorSetup{theorem=teal}
\section{Elegant Style Math Env}
\begin{axiom}[prime number]
  \boomen
\end{axiom}

\zlatexThmStyle{elegant}
\begin{definition}[prime number]
  \boomen
\end{definition}

\zlatexThmStyle{leftbar}
\begin{theorem}[prime number]
  \boomen
\end{theorem}

\zlatexThmTitleFormat{\textcolor{red}{\thmname{#1}:\thmnote{[#3]-}\thmnumber{#2}}}
\begin{lemma}[prime number]
  \boomen
\end{lemma}

\zlatexThmStyle{plain}
\begin{corollary}[prime number]
  \boomen
\end{corollary}

\newpage
\zlatexThmStyle{fancy}
\zlatexThmHook{begin=BEGIN, end=END, after=AFTER}
% \zlatexThmHook{begin=\begin{center}, end=\end{center}}
% \zlatexThmHook{before=\begin{framed}, after=\end{framed}}
\begin{proposition}
  \boomen
\end{proposition}


\zlatexThmStyleNew{before=\color{blue}\begin{framed}, after=\end{framed}}
\begin{definition}
  \boomen
\end{definition}

% \section{Final}
% Other trial as below:
% \begin{minted}{latex}
% \zslideThemeUse[
%     % UR={text=\FirstMark{2e-right-nonempty}} % failed
%     % UR={text=\FirstMark{test}}              % successed
%     % UR={text=\FirstMark{zslide-right}}      % successed
%     UR={text=\FirstMark{zslide-right}\, --\, \LastMark{zslide-right}} % successed
% ]{AnnArborDefault}
% % hyper link test
% \ExplSyntaxOn\makeatletter
% \int_step_inline:nnn {1}{4}{
%   % \hyper@link{link}{page.#1}{(#1)~|~}
%   \hyper@link{link}{zslide@HELLO.#1}{(HELLO.#1)~|~}
% }
% \makeatother\ExplSyntaxOff
% % navigate test
% \ExplSyntaxOn\makeatletter
% \cs_set:Npn \zslide@hyper@navigate:nn #1#2 {
%   % #1: total frame
%   % #2: current frame
%   \int_step_inline:nnn {1}{#1}{
%     \int_compare:nNnTF {#2} = {##1}
%       {\hyper@link{link}{zslide@HELLO.##1}{\textbullet}}
%       {\hyper@link{link}{zslide@HELLO.##1}{$\circ$}}
%   }
% }
% \zslide@hyper@navigate:nn {4}{2}
% \makeatother\ExplSyntaxOff
% \end{minted}

% \section{graphics}
% Auto scale graphics:

% \includegraphics[width=.2\linewidth]{Epmmy.jpg}
\end{document}









% ".color_set:nn" in l3keys
\documentclass{article}
\usepackage{xcolor}


\begin{document}
\ExplSyntaxOn\makeatletter
\cs_set:Npn \__temp__:nn #1#2 {\typeout{#1:#2}}
\cs_set:Npn \__temp__:n #1 {\typeout{#1}}
\keyval_parse:nnn 
  {\__temp__:n}
  {\__temp__:nn }
  {begin=AAA, end=BBB}


\def\ctext#1{\par\noindent\textcolor{\tl_use:c {l_test_#1_tl}}{Sample-TEXT}}


\cs_new:Npn \__zlatex_color_name_selector:w #1#2#3#4\q_stop 
  {#2}
\cs_new:Npn \__zlatex_color_init:nn #1#2 {
  \definecolor{zlatex@color@#1}{#1}{#2}\q_stop
  zlatex@color@#1
}

\tl_new:N \zlatex@color@keyC
\tl_set:Nn \zlatex@color@keyC {blue}
\keys_define:nn {test}{
  keyA  .tl_set:N   = \l_test_keyA_tl,
  keyA  .initial:n  = {red},
  keyB  .tl_set:N   = \l_test_keyB_tl,
  % keyB  .initial:n  = {blue},
  keyB  .initial:e  = {\__zlatex_color_name_selector:w \definecolor{orange}{HTML}{5f3dc4}\q_stop},
}
\ctext{keyA}
\ctext{keyB}
% \keys_set:nn {test}{keyA=black}
% \tl_show:N \l_test_keyB_tl

\par\noindent
\__zlatex_color_name_selector:w \definecolor{zlatex@color@keyC}{HTML}{5f3dc4}\q_stop

\makeatother\ExplSyntaxOff
\end{document}






\InputIfFileExists{zlatex-cfg.tex}{}{}
\documentclass[fancy, ]{../code/zlatex}
\zlatexSetup{mathSpec={envStyle=shadow}}
\zlatexThmCreate{theorem}{Ztheorem=THM|{HTML}{8e44ad}, Zproposition} 
% \zlatexColorSetup{
%   definition=blue,
%   theorem={HTML}{007175},
%   Ztheorem={RGB}{219,48,122}
% }
% \newsavebox{\zlatexBox}
% \zlatexThmStyleNew*{begin=\color{red}\begin{framed}, end=\end{framed}}

\def\boomen{As any dedicated reader can clearly see, the Ideal of practical
reason is a representation of, as far as I know, the things in themselves; 
\begin{align}
\underset{}{\mathbf{v} \bigotimes \mathbf{w}} 
  & = \underset{}{\mathbf{v} \otimes \mathbf{w}}
      = \sum_{i=1}^3\sum_{j=1}^3a_{ij}u^iv^j \\[-.75em]
  & = \sum_{i=1}^3\left(a_{i1}u^iv^1+a_{i2}u^iv^2+a_{i3}u^iv^3\right) 
  \end{align}  
}


\begin{document}
\section{FIRST}
\begin{definition}[DEF]
  Hello world: $a^2+b^2=c^2$.
\end{definition}

\zlatexThmStyle{elegant}
\begin{theorem}[theorem Test]
  Hello world: $a^2+b^2=c^2$.
\end{theorem}


\section{SECOND}
% \zlatexColorSetup{
%   % Ztheorem={HTML}{8e44ad},
%   Ztheorem=red 
% }
\begin{Ztheorem}[Ztheorem Test]
  Hello world: $a^2+b^2=c^2$.
\end{Ztheorem}

\begin{Zproposition}
  Hello world: $a^2+b^2=c^2$.
\end{Zproposition}


\newpage
% \zlatexColorSetup{theorem=teal}
\section{Elegant Style Math Env}
\begin{axiom}[prime number]
  \boomen
\end{axiom}

% \zlatexThmStyle{HOOK}
\begin{definition}[prime number]
  \boomen
\end{definition}

\zlatexThmStyle{leftbar}
\begin{theorem}[prime number]
  \boomen
\end{theorem}

\zlatexThmTitleFormat{\textcolor{red}{\thmname{#1}:\thmnote{[#3]-}\thmnumber{#2}}}
\begin{lemma}[prime number]
  \boomen
\end{lemma}

\zlatexThmStyle{plain}
\begin{corollary}[prime number]
  \boomen
\end{corollary}

\newpage
\zlatexThmStyle{fancy}
\zlatexThmHook{begin=BEGIN, end=END, after=AFTER}
% \zlatexThmHook{begin=\begin{center}, end=\end{center}}
% \zlatexThmHook{before=\begin{framed}, after=\end{framed}}
\begin{proposition}
  \boomen
\end{proposition}


\zlatexThmStyleNew{before=\color{blue}\begin{framed}, after=\end{framed}}
\begin{definition}
  \boomen
\end{definition}
\end{document}














\documentclass{article}
\usepackage{xcolor}

\begin{document}
\ExplSyntaxOn\makeatletter
\regex_new:N \l__color_mode_regex
\regex_set:Nn \l__color_mode_regex {(\cB..{1,}\cE.){2}}
\cs_new:Npn \test__color_set:n #1 {
  \regex_match:NnTF \l__color_mode_regex {#1}{
    \definecolor{Test@\l_keys_key_str}#1
    \tl_set:ce {l_test_\l_keys_key_str _tl}{Test@\l_keys_key_str}
  }{
    \@ifundefined{\string\color@#1}{
      \msg_new:nnn {color} {undefined} {--->~Color~`#1'~undefined}
      \msg_error:nn {color} {undefined}
    }{
      \tl_set:ce {l_test_\l_keys_key_str _tl}{#1}
    }
  }
}
\keys_define:nn {test}{
  keyA  .tl_set:N  = \l_test_keyA_tl,
  keyA  .initial:n = {green},
  keyA  .code:n    = {\test__color_set:n {#1}},
  keyB  .tl_set:N  = \l_test_keyB_tl,
  keyB  .initial:n = {black},
  keyB  .code:n    = {\test__color_set:n {#1}},
}

% EXAMPLES
\textcolor{\l_test_keyA_tl}{TEXT-keyA}\par
\textcolor{\l_test_keyB_tl}{TEXT-keyB}\par

\keys_set:nn {test}{keyA={HTML}{0f45f3}}
\textcolor{\l_test_keyA_tl}{TEXT-keyA}\par

\keys_set:nn {test}{keyA={HTML}{8975f3}}
\textcolor{\l_test_keyA_tl}{TEXT-keyA}\par

\keys_set:nn {test}{keyB=orange}
\textcolor{\l_test_keyB_tl}{TEXT-keyB}\par

\keys_set:nn {test}{keyB=orangee}
\textcolor{\l_test_keyB_tl}{TEXT-keyA}\par
\makeatother\ExplSyntaxOff
\end{document}






% hyperref test
\documentclass{article}
\usepackage[T1]{fontenc}
\usepackage{amsthm}
\newcounter{theorem}
\newenvironment{theorem}{\begin{flushleft}\textbf{Theorem}\;\refstepcounter{theorem}\thetheorem\enspace}{\end{flushleft}}
\newtheorem{definition}{Definition}
\makeatletter
\setlength{\parindent}{0pt}
\def\hyper@nopatch@thm{}
\usepackage{lipsum}
\usepackage{xcolor}
\usepackage{hyperref}
% \let\Hy@theorem@refstepcounter\stepcounter
\hypersetup{
  bookmarksnumbered,
  colorlinks = true,
  linkcolor = red,
  urlcolor = blue,
}
\def\seeTargets{%
  \par\noindent\dotfill\\
  \string\@currentHref\ = \meaning\@currentHref\par
  \string\@currentHpage\ = \meaning\@currentHpage\par
  \string\Hy@currentbookmarklevel\ = \meaning\Hy@currentbookmarklevel\par\vskip3em
}

\begin{document}
\tableofcontents

% \@ifundefined{hyper@nopatch@thm}{UN-DEF}{DEF}
% \meaning\@thm


\newpage
\section{Test 1}\label{test1}
\subsection{sss b}
Hello world I. 

AAAA=BBBB
% \MakeLinkTarget{ztargetA}
\hyper@anchor{ztargetB}
\seeTargets

\newpage
\section{thm link}
\begin{theorem}\label{thm1}
This is a theorem: $a^2+b^2=c^2$.
\end{theorem}

\begin{definition}\label{def1}
This is a definition: $a^2+b^2=c^2$.
\end{definition}

\seeTargets


\newpage
\section*{Test 2}
Hello world II, the star version section target is:\par
\seeTargets

\MakeLinkTarget[newType]{}
\seeTargets


\section{see targets}
\subsection{sss aaa}
\subsection{sss bbb}
\seeTargets


\section{Manually change target}
\subsection{theHCOUNTER}
Original current target is:\par
\seeTargets

Then i change current target manually:\par 
\def\theHsection{1000}
\seeTargets

\textbf{Can NOT change the target manually by \string\theHCOUNTER.}

\subsection{MakeLinkTarget-star}
The command \string\MakeLinkTarget* can change \string\@currentHref{} manually:\par
\MakeLinkTarget*{newManualType} \label{checkTarget}
\seeTargets

% \vskip3em
% \meaning\ref

% \vskip2em
% \ExplSyntaxOn
% \cs_meaning:N \__cmd_start:nNNnnn
% \cs_meaning:N \__cmd_start_aux:NNnnnn
% \cs_meaning:N \__cmd_run_code:
% \ExplSyntaxOff


\newpage
\section{book mark}
% \meaning\pdfbookmark\par
% \meaning\new@pdflink\par

AAA:\pdfbookmark[1]{ManuallyBookMark-I}{targetA}
\seeTargets

BBB: \currentpdfbookmark{ManuallyBookMark-II}{targetB}
\seeTargets

CCC: \subpdfbookmark{ManuallyBookMark-III}{targetC}
\seeTargets

DDD:(Even bookmark level not step, the number in dstination show step) \belowpdfbookmark{ManuallyBookMark-IV}{targetD}
\seeTargets

EEE: \hyper@anchorstart{target.X}\hyper@anchorend
\seeTargets


\newpage
1. \string\ref{} commad Test:
\begin{itemize}
  \item Hello world III:\ref{test1}
  \item Hello THM: \ref{thm1}
  \item Hello DEF: \ref{def1}
\end{itemize}

\vskip3em
2. \string\hyperlink{} commad Test:
\begin{itemize}
  \item Link Target Manually: % \LinkTargetOn ztargetA \LinkTargetOff
  \item \hyper@link {url}{ztargetB}{link text}; 
        \hyper@linkstart{link}{ztargetB} long long link text \hyper@linkend
\end{itemize}

\vskip3em
3. \string\MakeLinkTarget{} commad Test:
\begin{itemize}
  \item \hyper@link {link}{section*.2}{link text};
  \item \hyper@link {link}{newType*.3}{link text};
\end{itemize} 

\vskip3em
4. bookmark related target link:
\begin{itemize}
  \item AAA: \hyper@link {link}{targetA.1}{link text 1};
  \item BBB: \hyper@link {link}{targetB.1}{link text 2};
  \item CCC: \hyper@link {link}{targetC.2}{link text 3};
  \item DDD: \hyper@link {link}{targetD.3}{link text 4};
  \item EEE: \hyper@link {link}{target.X}{link text 5};
\end{itemize} 
\end{document}





\InputIfFileExists{zlatex-cfg.tex}{}{}
\documentclass[
  fancy,
  % hyper,
  % lang=cn,
  % class=book,
  % font={config},
  % layout={margin},
  % layout={slide, theme=AnnArborSeahorse, aspect=16|9},
  % classOption={12pt}
]{../code/zlatex}
\zlatexThmLang{fr}
\zlatexloadlibrary{mathalias}
\zlatexSetup{mathSpec={envStyle=fancy}}
% \zlatexColorSetup{theorem=teal}
\zlatexThmCnt{share, parent=section}
\zlatexThmCreate{theorem}{Zaxiom, Ztheorem=Thm|green, Zproposition=Prop|orange} %Ztheorem=Thm|{HTML}{0f45f3}
\zlatexThmCreate{proof}{Zproof, Zexample=Example|red, Zsolution=Solution|}


\def\boomen{As any dedicated reader can clearly see, the Ideal of practical
reason is a representation of, as far as I know, the things in themselves; 
\begin{align}
\underset{}{\mathbf{v} \bigotimes \mathbf{w}} 
  & = \underset{}{\mathbf{v} \otimes \mathbf{w}}
      = \sum_{i=1}^3\sum_{j=1}^3a_{ij}u^iv^j \\[-.75em]
  & = \sum_{i=1}^3\left(a_{i1}u^iv^1+a_{i2}u^iv^2+a_{i3}u^iv^3\right) 
  \end{align}  
}

% \setlength{\fboxsep}{0pt}
\title{z\LaTeX{} Benchmark Test}
\author{Eureka}
\date{\today}
\begin{document}
\maketitle
% \chapter{Hello}
% \section{Py Fijw}
\section{FIRST}
Hello world; \meaning\FF


\begin{theorem}[Pythagorean theorem]\label{Pythagorean theorem}
This is a theorem:
  \begin{align}
    \R{d}^2x + \R{d}^2y = \R{d}^2z\\
    dx = 2
  \end{align}  
\end{theorem}

\begin{proof}\label{proof-1}
This is a proof: $a^2 + b^2 = c^2$
\end{proof}


\section{Internal Math Env}
\begin{zlatexZtheorem}[internal Env]
  Hello Internal Ztheorem
\end{zlatexZtheorem}


\section{New Math Env}
\begin{Zaxiom}
Hello Zaxiom
\end{Zaxiom}

\begin{Ztheorem}[New Thm Users' ENV] % ---> inside thm source
\boomen

\boomen
\end{Ztheorem}

\zlatexThmStyle{elegant}
\begin{Zproposition}[New Prop Users' ENV]\label{zprop-1}
\boomen

\boomen
\end{Zproposition}



\newpage\zlatexColorSetup{theorem=teal}
\section{Elegant Style Math Env}
\begin{axiom}[prime number]
  \boomen
\end{axiom}

\begin{definition}[prime number]
  \boomen
\end{definition}

\begin{theorem}[prime number]
  \boomen
\end{theorem}

\zlatexThmTitleFormat{\thmname{#1}:\thmnote{[#3]-}\thmnumber{#2}}
\begin{lemma}[prime number]
  \boomen
\end{lemma}

\zlatexThmStyle{plain}
\begin{corollary}[prime number]
  \boomen
\end{corollary}

\begin{proposition}
  \boomen
\end{proposition}


\section{New Proof Env}
\begin{Zproof}
Hello Zproof
\end{Zproof}

\begin{Zexample}
Hello Zexample
\end{Zexample}

\begin{Zsolution}\label{zsolu-1}
Hello Zsolution
\end{Zsolution}



% \href{http://www.google.com}{Google}

% \newlength{\myl}
% \settowidth\myl{\hbox{\Large\bfseries Eureka-Eureka-Eureka-Eureka-Eureka-}}
% \the\myl % 289.98325pt


% \ExplSyntaxOn
% \__zlatex_math_env_color_set:w thm|red\q_stop
% <==>
% \exp_last_unbraced:Ne \__zlatex_math_env_color_set:w {thm|red}\q_stop
% \tl_use:c {l__zlatex_thm_color_tl}
% \ExplSyntaxOff

\newpage
Hello world:

\begin{itemize}
  \item \cref{Pythagorean theorem}
  \item \cref{proof-1}
  \item \cref{zprop-1}
  \item \cref{zsolu-1}
\end{itemize}


% \includegraphics[width=.5\paperwidth]{example-image-a}
\end{document}

















\documentclass[
  % fancy,
  % hyper,
  % lang=cn,
  % class=book,
  % classOption={12pt}
  % font={config},
  % layout={margin},
  % layout={slide, theme=AnnArborSeahorse, aspect=16|9},
]{../code/zlatex}
\zlatexloadlibrary{mathalias}
% \zlatexThmCnt{share, parent=section}


\begin{document}
Hello world: $\int x \dd x$.
\begin{align}
  & \int x \dd x = \frac{1}{2}x^2 + \CC \\
  & e^{\int x \dd x} = e^{\frac{1}{2}x^2 + \CC} \\
  & \K{A} = \S{A} = \C{A} = \B{A} \\
  & A\pi+\sum = \F{A\pi+\sum} 
\end{align}


\begin{align}
  & \grad \cdot F = \div F \\
  & \grad F = \div F \\
  & \olddiv F \sim \oldhom F \\
  & \cong = \backsimeq
\end{align}


\begin{align}
  & \A a \xl b \xl[\alpha]((\psi)) c \\
  & a \ma b \xxl(\alpha) c \\
  & \E a \xr(\Theta = \Gamma + \zeta) b \xr[\alpha] c \\
  & a \da b \xxr{[}{[}\alpha{]}{]}(\beta - \gamma = \xi) c 
\end{align}
\end{document}







\documentclass{article}
\usepackage{amsmath, amsthm}


% preamble
\ExplSyntaxOn
\NewHook{zlatex/thm/titleformat}
\newtheoremstyle{zlatexMathEnv}
  {2pt}{2pt}{}
  {0pt}{\bfseries}{}
  {.25em}
  {
    \cs_set:Npn \ThmName { \thmname{#1} }
    \cs_set:Npn \ThmNumber { \thmnumber{#2} }
    \cs_set:Npn \ThmNote { \thmnote{#3} }
    % \tl_if_empty:nTF {#3}
    %   { \ThmName \ThmNumber }
    %   {  }
    \UseHook{zlatex/thm/titleformat}
  }
\theoremstyle{zlatexMathEnv}

\newcommand{\zlatexThmTitleFormat}[1]{
  \IfHookEmptyTF{zlatex/thm/titleformat}{}
    {\RemoveFromHook{zlatex/thm/titleformat}[once]}
  \AddToHook{zlatex/thm/titleformat}[once]{#1}
}
\ExplSyntaxOff


% main document
\begin{document}
\newtheorem{aaa}{AAA}
\zlatexThmTitleFormat{\ThmName. \ThmNumber [\ThmNote)}
\begin{aaa}[SECOND]
  ENV-CONTENT-SECOND
\end{aaa}
\end{document}









\documentclass{article}
\usepackage{amsmath, amsthm}


\begin{document}
\ExplSyntaxOn
\NewHookWithArguments{zlatex/thm/titleformat}{3}
\newtheoremstyle{zlatexMathEnv}
  {2pt}{2pt}{}
  {0pt}{\bfseries}{}
  {.25em}
  {
    \cs_set:Npn \ThmName { \thmname{#1} }
    \cs_set:Npn \ThmNumber { \thmnumber{#2} }
    \cs_set:Npn \ThmNote { \thmnote{#3} }
    % { \ThmName \ThmNumber \ThmNote }
    \UseHookWithArguments{zlatex/thm/titleformat}{3}{#1}{#2}{#3}
  }
\theoremstyle{zlatexMathEnv}

% \tl_put_right:Nn \l_text_expand_exclude_tl {\thmnote \thmnumber \thmname}
% \cs_generate_variant:Nn \text_expand:n {e}
% \newcommand{\foo}[1]{
%   \tl_set:Nn \l_tmpa_tl {\text_expand:n {#1}}
%   \tl_show:N \l_tmpa_tl
% }
% \foo{\thmname{##1} \thmnumber{#2} \thmnote{[#3]}}


% ===>
% \cs_new_protected:Nn \keythms_thmstyle_thmname:n { \thmname{#1} }
% \cs_new_protected:Nn \keythms_thmstyle_thmnumber:n { \thmnumber{#1} }
% \cs_new_protected:Nn \keythms_thmstyle_thmnote:n { \thmnote{#1} }

% \tl_put_right:Nn \l_text_expand_exclude_tl { \thmnote \thmnumber \thmname }
% % ^ allows \thmnote, etc. to work in headstyle; hope no bad side effects!
% % user is in charge of spacing with \NAME and \NUMBER (thmtools compat...)
% \cs_set:Npn \NAME { \keythms_thmstyle_thmname:n { ##1 } }
% \cs_set:Npn \NUMBER
%   {
%     \keythms_thmstyle_thmnumber:n
%       {
%         \group_begin:
%         ##2
%         \group_end:
%       }
%   }
% \cs_set:Npn \NOTE
%   {
%     \keythms_thmstyle_thmnote:n {##3}
%   }

% \text_expand:n { \keythms_thmstyle_headcmd:nnn{##1}{##2}{##3} }
% \cs_set:Nn \keythms_thmstyle_headcmd:nnn
% {
%   \makebox[0pt][r]{
%     \keythms_thmstyle_thmnumber:n {
%         ##2
%       }
%     }
%   \NAME\NOTE
% }
% \foo{\NAME \NUMBER \NOTE}
% \foo{\thmname{##1} \thmnumber{##2} \thmnote{[##3]}}



\newcommand{\zlatexThmTitleFormat}[1]{
  \IfHookEmptyTF{zlatex/thm/titleformat}{}
    {\RemoveFromHook{zlatex/thm/titleformat}[once]}
  \AddToHookWithArguments{zlatex/thm/titleformat}[once]{#1}
}
\ExplSyntaxOff
\newtheorem{aaa}{AAA}

% \zlatexThmTitleFormat{\thmname{#1}: \thmnumber{#2} \thmnote{[#3]}}
\begin{aaa}[FIRST]
  ENV-CONTENT-FIRST
\end{aaa}

% \zlatexThmTitleFormat{\thmname{#1}. \thmnumber{#2} \thmnote{(#3)}}
\zlatexThmTitleFormat{\ThmName. \ThmNumber [\ThmNote)}
\begin{aaa}[SECOND]
  ENV-CONTENT-SECOND
\end{aaa}

\begin{aaa}
  ENV-CONTENT-THIRD
\end{aaa}
\end{document}








\documentclass[
  layout={
    slide, 
    aspect=16|9, 
    % theme=AnnArborAlbatross
  },
  % mathSpec={
  %   font=hello,
  %   world=haha
  % },
  % lang=cn,
  packageOption={
    % geometry={margin=3in},
    % fontspec={quiet},
    % fontenc={T1}
  }
]{../code/zlatex}
% \geometry{margin=3in}%{geometry}
% \documentclass[
%   fancy,
%   hyper,
%   % class=hello,
%   classOption={12pt},
%   layout={margin, slide, aspect=16|9, theme=AnnArborAlbatrosss},
%   % font=config,
%   toc={titlee=CONTENTS},
%   mathSpec={
%     font=eulerr, 
%     counter=world
%   },
%   font=lm,
%   layout=openAll,
%   bib_index={Notload},
%   % lang=cn,
%   hello=world,
% ]{../code/zlatex}
% \usepackage{pifont}
% \usepackage{anyfontsize}
\zlatexSetup{mathSpec={envStyle=fancy, font=mathpazo}}
% \zslideSetup{
%   world=haha,
%   toc={hello=hi}
% }



% \documentclass{article}


\title{Debug}
\author{Eureka}
\date{\today}
\begin{document}
\maketitle
% \tableofcontents
% \newpage


Hello world
\begin{theorem}
This is Hello world:

\begin{align}
  \sum_{i=1}^{+\infty}{\int_{0}^{i}-\frac{1}{t}\mathrm{d}t} = \frac{\pi^2}{6}
\end{align}
\end{theorem}


% \newpage 
% \section{Page 2}
% \subsection{sss}
% world

% \newpage
% \section{Math Formula}
% \begin{align}
%   \sum_{i=1}^{+\infty}{\int_{0}^{i}-\frac{1}{t}\mathrm{d}t} = \frac{\pi^2}{6}
% \end{align}


\ExplSyntaxOn
\ExplSyntaxOff

\end{document}
